\chapter*{ВВЕДЕНИЕ}
\addcontentsline{toc}{chapter}{ВВЕДЕНИЕ}
\sloppy

В последние годы наблюдается устойчивый рост интереса к агентным системам, в частности к многоагентным архитектурам на базе больших языковых моделей. Это связано с переходом от статических моделей обработки данных к более гибким системам, способным автономно планировать, декомпозировать задачи и взаимодействовать с внешними инструментами. Особенно заметен рост числа работ, посвященных применению таких систем в исследовательских задачах: автоматизация обзоров литературы, генерация гипотез, постановка и проведение вычислительных экспериментов. Однако, несмотря на активное развитие, большинство существующих решений (включая архитектуры, аналогичные OpenClaw и другим агентным фреймворкам) остаются ограниченными с точки зрения устойчивости, управляемости и воспроизводимости. Сейчас такие системы демонстрируют нестабильное качество при решении сложных исследовательских задач, что формирует необходимость в разработке архитектурных улучшений, направленных на повышение их эффективности, интерпретируемости и воспроизводимости.

Практическая значимость данной работы заключается в разработке улучшенной архитектуры агентной системы, ориентированной на автоматизацию научного цикла. Предлагаемые изменения позволяют повысить качество работы системы в задачах анализа, генерации и проверки гипотез, что может быть применимо в широком спектре областей - от анализа научных публикаций до поддержки принятия решений в R&D. Улучшение механизмов памяти и управления агентами потенциально снижает затраты времени и ресурсов при выполнении сложных интеллектуальных задач. Теоретическая значимость работы связана с вкладом в развитие архитектурных подходов к построению агентных систем. В частности, рассматривается переход от линейных к структурированным представлениям знаний, формализация цикличности научного процесса в агентной логике, а также постановка задачи динамической оптимизации выбора моделей. Эти аспекты расширяют существующие представления о том, как должны быть организованы когнитивные процессы в системах, основанных на языковых моделях, и могут служить основой для дальнейших исследований в области интеллектуальных агентов.

Целью данной работы является разработка и экспериментальная оценка архитектурных улучшений агентной системы на базе OpenClaw и ClawdLab, направленных на повышение ее эффективности при решении исследовательских задач. 

Для достижения поставленной цели формулируются следующие задачи: 
\begin{itemize}
  \item проанализировать существующие архитектуры агентных систем и выявить их ограничения, в том числе в контексте решения научных задач. 

  \item разработать механизм структурированной памяти, позволяющий извлекать  и хранить информацию в структурном виде для эффективной работы над задачами с большими контекстами

  \item разработать механизм цикличности, обеспечивающий более гибкое выполнение этапов научного процесса

  \item реализовать модуль оптимизации выбора наиболее подходящих больших языковых моделей, в зависимости от текущей подзадачи

  \item интегрировать предложенные компоненты в единую архитектуру

  \item провести экспериментальную оценку разработанной системы и сравнить ее с базовой архитектурой по заданным критериям
\end{itemize}

Новизна работы носит прикладной характер и заключается в модификации существующей архитектуры агентной системы с целью устранения выявленных ограничений. В отличие от традиционного подхода, где память агента представлена в виде последовательного лога, предлагается использовать структурированную память, вводится модель цикличности, отражающая природу научного процесса, а также внедряется модуль оптимизации выбора моделей, который динамически адаптирует поведение системы в зависимости от контекста задачи. В совокупности эти изменения формируют более гибкую и адаптивную архитектуру, способную лучше справляться со сложными многошаговыми задачами. Таким образом, работа вносит вклад в область инженерии агентных систем, предлагая конкретные архитектурные решения и демонстрируя их эффективность.

Для оценки эффективности предложенных архитектурных улучшений используется экспериментальный подход, основанный на сравнении базовой и модифицированной систем. Эксперименты проводятся на наборе задач, имитирующих элементы научного цикла (например, анализ текста, генерация гипотез, выбор методов и интерпретация результатов). 

В качестве критериев оценки выделяются три ключевых аспекта: 
\begin{itemize}
  \item качество решений, которое может измеряться через метрики точности, полноты или экспертной оценки результатов
  \item устойчивость и воспроизводимость, отражающая способность системы получать согласованные результаты при повторных запусках и варьировании входных данных
  \item эффективность, включая количество шагов, время выполнения и использование ресурсов. 
\end{itemize}

Сравнение проводится в контролируемых условиях с фиксированными сценариями, что позволяет изолировать влияние предложенных изменений. Успехом работы считается значимое улучшение по указанным критериям по сравнению с базовой архитектурой, а также демонстрация более предсказуемого и интерпретируемого поведения системы.
